% ── Korelační analýza: pokrytí KS a ukazatele trhu práce ─────────────────────

\paragraph{Datová základna a pokrytí zemí EU27}\label{par:korelace_data}

Korelační analýza pracuje s~párovým panelem (\textit{geo\,$\times$\,rok}) sestaveným z~následujících zdrojů.
Pokrytí \ac{KS} pochází z~databáze OECD/AIAS ICTWSS~\cite{oecd_aias_ictwss_CBC_ERB_pct}: pro 21~zemí EU27 je použita upravená míra \textit{AdjCov} zachycující podíl zaměstnanců pokrytých jakoukoliv kolektivní smlouvou bez ohledu na odborové členství; pro Německo~(DE), Slovensko~(SK) a~Slovinsko~(SI) je \textit{AdjCov} nahrazena mírou ERB z~průzkumu OECD/CBC, neboť \textit{AdjCov} pro tyto země stagnuje přibližně na hodnotách z~roku 1990 resp.~2015--2016.
Hustota odborů pochází z~téže databáze (ukazatel TUD~\cite{oecd_aias_ictwss_TUD_pct}).
Výstupní proměnné pocházejí z~Eurostatu: průměrná skutečná týdenní pracovní doba (lfsa\_ewhan2~\cite{eurostat_lfsa_ewhan2_HR_weekly}), gender pay gap v~nefinančním sektoru (earn\_gr\_gpgr2~\cite{eurostat_gpg}), čistý hodinový příjem v~PPS odvozený podílem ročního příjmu a~skutečné roční pracovní doby (earn\_nt\_net~\cite{eurostat_earn_nt_net_PPS_AW100}, lfsa\_ewhan2), produktivita práce na hodinu v~PPS (nama\_10\_lp\_ulc~\cite{eurostat_nama_10_lp_ulc_NLPR_HW_EU27eq100}) a~Giniho koeficient disponibilního příjmu (ilc\_di12~\cite{eurostat_ilc_di12_Gini}).
Datová řada začíná rokem 2004 shodně ve všech souborech.

\paragraph{Odvození hodinového čistého příjmu v PPS}\label{par:pps_hourly}

Eurostatový soubor earn\_nt\_net~\cite{eurostat_earn_nt_net_PPS_AW100}
poskytuje roční čistý příjem standardizované osoby ve výši
\SI{100}{\percent} průměrné mzdy (\textit{average worker}, AW100) v~jednotkách
\SI{}{\pps}, tj.~$\acs{var-y}^{\acs{idx-PPS}}_{\acs{idx-n},\acs{idx-AW100},\acs{idx-a}}$.
Pro srovnání s~hodinovou produktivitou je nutný přepočet na hodinovou sazbu.
Roční pracovní fond~$\acs{var-H}_{\acs{idx-a}}$ se odhaduje z~průměrné skutečné
týdenní pracovní doby~$\acs{var-H}_{\acs{idx-w}}$ (Eurostat
lfsa\_ewhan2~\cite{eurostat_lfsa_ewhan2_HR_weekly}) přepočtem na 52,18~týdne
v~roce (průměrná délka roku $365{,}2425/7 \approx 52{,}18$):
\begin{equation}\label{eq:pps_hourly}
  \acs{var-w}^{\acs{idx-PPS}}_{\acs{idx-n},\acs{idx-h},\acs{idx-AW100}}
    = \frac{\acs{var-y}^{\acs{idx-PPS}}_{\acs{idx-n},\acs{idx-AW100},\acs{idx-a}}}
           {\acs{var-H}_{\acs{idx-w}} \cdot 52{,}18}.
\end{equation}
Výsledná veličina~$\acs{var-w}^{\acs{idx-PPS}}_{\acs{idx-n},\acs{idx-h},\acs{idx-AW100}}$
[\SI{}{\pps\per\hour}] zohledňuje jak rozdíly v~čisté koupěschopnosti, tak
v~délce skutečného pracovního fondu, a~je proto srovnatelnější napříč zeměmi
než hrubá hodinová mzda nebo roční čistý příjem samostatně.

Ze 27~členských zemí EU jsou do analýzy zahrnuty \textbf{čtyřiadvacet}; z~dat jsou vyjmuty tři:
\begin{itemize}
  \item \textbf{Irsko~(IE)} je záměrně vyřazeno jako statistický outlier.
    Jeho HDP/ob.~v~PPS dosahuje přibližně 237~(EU27\,=\,100), přičemž tato hodnota je způsobena přesunem zisků nadnárodních technologických a~farmaceutických korporací se sídlem v~Irsku pro účely daňové optimalizace v~prostoru EU (transfer pricing, relokace duševního vlastnictví).
    Skutečná životní úroveň, vyjádřená upravenými národními příjmy GNI* (CSO Ireland), se pohybuje okolo indexu~135.
    Zahrnutí Irska by systematicky zkreslilo regresní koeficient $\hat{\beta}$ u~páru pokrytí~vs.~produktivita a~tím deformovalo predikci pro zbývající země.
  \item \textbf{Estonsko~(EE) a~Lotyšsko~(LV)} --- databáze ICTWSS \textit{AdjCov} pro tyto země neobsahuje žádné záznamy; jde o~strukturálně chybějící hodnoty (MCAR), jejichž doplnění bez externího zdroje není opodstatněné.
  \item \textbf{Lucembursko~(LU)} --- data \textit{AdjCov} jsou dostupná pouze před rokem 2004; po aplikaci filtru \texttt{START\_YEAR\,=\,2004} vypadávají z~panelu.
\end{itemize}

Zbývajících 24~zemí pokrývá všechny geografické a~institucionální skupiny EU27: nordické, kontinentální, středomořské, středoevropské a~balkánské (s~výjimkou EE a~LV).
Panel není vyvážený: počet párových pozorování~$\acs{var-n}_{\text{panel}}$ se liší dle dostupnosti výstupní proměnné (viz tabulku~\ref{tab:korelace_tabulka}).

\paragraph{Matematický aparát}\label{par:korelace_math}

Jsou použity tři statistické míry asociace a~jeden predikční model.

\emph{Pearsonův korelační koeficient} měří sílu a~směr lineárního vztahu mezi proměnnými~$\acs{var-x}$ a~$\acs{var-y}$:
\begin{equation}\label{eq:pearson_r}
  \acs{var-r}_{\acs{idx-xy}}
  = \frac{\displaystyle\sum_{\acs{idx-i}=1}^{\acs{var-n}}
            (\acs{var-x}_{\acs{idx-i}} - \bar{x})\,
            (\acs{var-y}_{\acs{idx-i}} - \bar{y})}
         {\displaystyle\sqrt{
            \sum_{\acs{idx-i}=1}^{\acs{var-n}}
              (\acs{var-x}_{\acs{idx-i}}-\bar{x})^2
            \cdot
            \sum_{\acs{idx-i}=1}^{\acs{var-n}}
              (\acs{var-y}_{\acs{idx-i}}-\bar{y})^2}}
  \;\in\; [-1,\,1].
\end{equation}
Tento vztah je využit v~modelovém výpočtu pro ČR (viz~rov.~\eqref{eq:model_cz}).
Síla korelací je hodnocena dle Evansovy (1996) škály: $|\acs{var-r}| < 0{,}20$ velmi slabá;
$0{,}20$--$0{,}39$ slabá; $0{,}40$--$0{,}59$ střední; $0{,}60$--$0{,}79$ silná;
$\geq 0{,}80$ velmi silná.

\emph{Spearmanův koeficient~$\acs{var-rho}$} je Pearsonovo~$\acs{var-r}$ aplikované na hodnosti (ranky) pozorování namísto původních hodnot:
\begin{equation}\label{eq:spearman_rho}
  \acs{var-rho} = \acs{var-r}\!\bigl(R(\acs{var-x}),\,R(\acs{var-y})\bigr)
       = 1 - \frac{6\,\displaystyle\sum_{\acs{idx-i}=1}^{\acs{var-n}} d_{\acs{idx-i}}^{\,2}}
              {\acs{var-n}\,(\acs{var-n}^{2}-1)},
  \quad d_{\acs{idx-i}} = R(\acs{var-x}_{\acs{idx-i}}) - R(\acs{var-y}_{\acs{idx-i}}),
\end{equation}
kde pravý výraz platí při absenci shodných hodnot (\textit{ties}); obecně platí levý tvar.
Výrazná divergence~$|\acs{var-r} - \acs{var-rho}|$ signalizuje přítomnost vlivných odlehlých pozorování, která posouvají lineární vztah.

\emph{Prostá lineární regrese} metodou nejmenších čtverců (OLS):
\begin{equation}\label{eq:ols}
  \acs{var-yhat} = \hat{\alpha} + \hat{\beta}\,\acs{var-x},
  \qquad
  \hat{\beta}
    = \frac{\displaystyle\sum_{\acs{idx-i}=1}^{\acs{var-n}}
              (\acs{var-x}_{\acs{idx-i}}-\bar{x})\,
              (\acs{var-y}_{\acs{idx-i}}-\bar{y})}
           {\displaystyle\sum_{\acs{idx-i}=1}^{\acs{var-n}}
              (\acs{var-x}_{\acs{idx-i}}-\bar{x})^{2}},
  \quad
  \hat{\alpha} = \bar{y} - \hat{\beta}\,\bar{x}.
\end{equation}
Parametr~$\hat{\beta}$ udává průměrnou změnu $\acs{var-yhat}$ při nárůstu pokrytí \ac{KS} o~jeden procentní bod.

\emph{Koeficient determinace}~$\acs{var-Rdet}$ vyjadřuje podíl rozptylu~$\acs{var-y}$ vysvětlený regresním modelem:
\begin{equation}\label{eq:r2}
  \acs{var-Rdet}
  = 1 - \frac{\displaystyle\sum_{\acs{idx-i}=1}^{\acs{var-n}}
                (\acs{var-y}_{\acs{idx-i}} - \acs{var-yhat}_{\acs{idx-i}})^{2}}
             {\displaystyle\sum_{\acs{idx-i}=1}^{\acs{var-n}}
                (\acs{var-y}_{\acs{idx-i}} - \bar{y})^{2}}
  = \acs{var-r}_{\acs{idx-xy}}^{\,2} \;\in\; [0,\,1],
\end{equation}
kde rovnost $\acs{var-Rdet} = \acs{var-r}_{\acs{idx-xy}}^{\,2}$ platí pro prostou regresi s~konstantou.

\emph{Predikce dopadu změny pokrytí}~$\Delta \acs{var-x}$ na výstupní proměnnou:
\begin{equation}\label{eq:delta_pred}
  \Delta\acs{var-yhat} = \hat{\beta}\,\Delta \acs{var-x}.
\end{equation}

\paragraph{Kombinovaný pohled na asociace pokrytí KS s trhem práce}\label{par:scatter_combined}

Kombinovaný rozptylový diagram (obr.~\ref{fig:korelace_scatter}) zachycuje simultánně vztah pokrytí \ac{KS} ke čtyřem výstupním proměnným a~umožňuje evaluaci omnibusové hypotézy práce: silnější \ac{KV} instituce jsou systematicky asociovány s~lepšími výsledky napříč všemi sledovanými dimenzemi.
Každý dílčí panel potvrzuje parciální asociaci, ale kombinovaný pohled odhaluje koherentní systémový vzorec: státy s~vysokým pokrytím \ac{KS} (AT, DK, BE, FR) konzistentně vykazují vyšší hodinové příjmy v~PPS, vyšší produktivitu práce, nižší pracovní dobu za srovnatelný příjem a~nižší gender pay gap.
ČR se ve všech čtyřech dílčích grafech nachází v~kvadrantu nízkého pokrytí, ačkoli absolutně není na dně EU27.
Polysystémový argument je silnější než jakákoli dílčí korelace: holistické posílení \ac{KV} pokrytí (mechanismus erga omnes, větší dopad \ac{APZ}) by s~vysokou pravděpodobností zlepšilo celou sadu výstupů současně, protože sdílejí společný institucionální determinant.

% Editable figure definition for korelace_scatter
% Edit freely — Python will NOT overwrite this file.
% To regenerate defaults, delete this file and re-run the script.
\def\strKorelaceScatterCaption{\centering Korelace pokrytí KV a veličin trhu práce, EU27. Zdroj dat: OECD\,ICTWSS~\cite{oecd_aias_ictwss_CBC_ERB_pct}, Eurostat~\cite{eurostat_lfsa_ewhan2_HR_weekly}\cite{eurostat_gpg}\cite{eurostat_earn_nt_net_PPS_AW100}\cite{eurostat_nama_10_lp_ulc_NLPR_HW_EU27eq100}.\protect\footnotemark}%
%
\begin{figure}[htbp]
  \centering
  \resizebox{0.98\linewidth}{!}{% Editable figure definition for korelace_scatter
% Edit freely — Python will NOT overwrite this file.
% To regenerate defaults, delete this file and re-run the script.
\def\strKorelaceScatterCaption{\centering Korelace pokrytí KV a veličin trhu práce, EU27. Zdroj dat: OECD\,ICTWSS~\cite{oecd_aias_ictwss_CBC_ERB_pct}, Eurostat~\cite{eurostat_lfsa_ewhan2_HR_weekly}\cite{eurostat_gpg}\cite{eurostat_earn_nt_net_PPS_AW100}\cite{eurostat_nama_10_lp_ulc_NLPR_HW_EU27eq100}.\protect\footnotemark}%
%
\begin{figure}[htbp]
  \centering
  \resizebox{0.98\linewidth}{!}{% Editable figure definition for korelace_scatter
% Edit freely — Python will NOT overwrite this file.
% To regenerate defaults, delete this file and re-run the script.
\def\strKorelaceScatterCaption{\centering Korelace pokrytí KV a veličin trhu práce, EU27. Zdroj dat: OECD\,ICTWSS~\cite{oecd_aias_ictwss_CBC_ERB_pct}, Eurostat~\cite{eurostat_lfsa_ewhan2_HR_weekly}\cite{eurostat_gpg}\cite{eurostat_earn_nt_net_PPS_AW100}\cite{eurostat_nama_10_lp_ulc_NLPR_HW_EU27eq100}.\protect\footnotemark}%
%
\begin{figure}[htbp]
  \centering
  \resizebox{0.98\linewidth}{!}{\input{../python/figures/korelace_scatter.pgf}}
  \caption{\strKorelaceScatterCaption}
  \label{fig:korelace_scatter}
\end{figure}
\footnotetext{Irsko (IE) vyřazeno: HDP/ob.~v~PPS $\approx 237$ (EU27\,=\,100) je zkresleno přesunem zisků nadnárodních korporací se sídlem v~Irsku (transfer pricing, relokace~IP) --- nereprezentuje skutečnou životní úroveň ani podmínky trhu práce. Data nedostupná: EE, LV --- chybějící hodnoty v~ICTWSS~\textit{AdjCov}; LU --- data dostupná pouze před rokem 2004.}
}
  \caption{\strKorelaceScatterCaption}
  \label{fig:korelace_scatter}
\end{figure}
\footnotetext{Irsko (IE) vyřazeno: HDP/ob.~v~PPS $\approx 237$ (EU27\,=\,100) je zkresleno přesunem zisků nadnárodních korporací se sídlem v~Irsku (transfer pricing, relokace~IP) --- nereprezentuje skutečnou životní úroveň ani podmínky trhu práce. Data nedostupná: EE, LV --- chybějící hodnoty v~ICTWSS~\textit{AdjCov}; LU --- data dostupná pouze před rokem 2004.}
}
  \caption{\strKorelaceScatterCaption}
  \label{fig:korelace_scatter}
\end{figure}
\footnotetext{Irsko (IE) vyřazeno: HDP/ob.~v~PPS $\approx 237$ (EU27\,=\,100) je zkresleno přesunem zisků nadnárodních korporací se sídlem v~Irsku (transfer pricing, relokace~IP) --- nereprezentuje skutečnou životní úroveň ani podmínky trhu práce. Data nedostupná: EE, LV --- chybějící hodnoty v~ICTWSS~\textit{AdjCov}; LU --- data dostupná pouze před rokem 2004.}


\paragraph{Tabulka korelačních koeficientů}\label{par:corr_table}

Tabulka~\ref{tab:korelace_tabulka} shrnuje výsledky korelační analýzy za roky 2004--2024.
Každý řádek reportuje dvě úrovně analýzy:
\begin{enumerate}
  \item \textbf{Panelový koeficient} (Pearsonovo~$\acs{var-r}$, Spearmanovo~$\acs{var-rho}$): počítán na všech dostupných párech geo$\times$rok ($\acs{var-n}_{\text{panel}}$).
    Tato metoda zachycuje jak \emph{mezi-zemní} (between), tak \emph{vnitro-zemní} (within) variaci --- tj.\ jak rozdíly \emph{mezi} zeměmi, tak změny \emph{v~čase v~rámci} jedné země.
    Pozorování však nejsou nezávislá: stejná země se v~panelu vyskytuje opakovaně, čímž vzniká autokorelace uvnitř klastru.
    Sloupec~$\acs{var-n}_{\text{geo}}$ udává počet odlišných zemí v~panelu.
  \item \textbf{Průřezový koeficient} ($\acs{var-r}_{2024}$): Pearsonovo~$\acs{var-r}$ pouze pro nejnovější společný rok ($\acs{var-n}_{2024}$~zemí).
    Jde o~klasický průřez (cross-section), kde každá země vstupuje právě jednou --- pozorování jsou proto nezávislá.
    Nižší $\acs{var-n}$ vede k~menší statistické síle, ale výsledek netrpí vnitro-klastrovým opakováním.
\end{enumerate}
Obě úrovně se doplňují: panelový koeficient má vyšší statistickou sílu díky většímu~$\acs{var-n}$, průřezový koeficient zachycuje aktuální stav bez historické zátěže.

Tabulka \ref{tab:korelace_tabulka} shrnuje výsledky korelační analýzy provedené na párovém panelu zemí OECD/\ac{EU} za dostupná roky od roku 2004.
Zdrojem dat pro pokrytí kolektivními smlouvami je databáze OECD/AIAS ICTWSS (upravená míra \textit{AdjCov} pro většinu zemí \acs{geo-EU27}; pro Německo a~Slovensko průzkumová míra ERB, viz popis dat); pro odborovou organizovanost tatáž databáze (TUD); výstupní proměnné pocházejí z~Eurostatu (lfsa\_ewhun2, nama\_10\_pc, ilc\_di12, earn\_nt\_net).
Z~analýzy je vyřazeno Irsko: jeho \ac{HDP}/ob.~v~\acs{PPS} $\approx 237$ (\acs{geo-EU27}\,=\,100) je zkresleno přesunem zisků nadnárodních korporací se sídlem v~Irsku (transfer pricing, relokace duševního vlastnictví) --- nereprezentuje skutečnou životní úroveň ani podmínky trhu práce.
Estonsko a~Lotyšsko chybí kvůli neúplným hodnotám pokrytí \acs{KS} v~ICTWSS~\textit{AdjCov}; pro Lucembursko jsou data pokrytí dostupná pouze před rokem 2004.
Každý řádek reportuje: počet párových pozorování $\acs{var-n}$ (geo$\times$rok po párování datových sad), Pearsonovo $\acs{var-r}$ na celém panelu (míra lineární asociace), Spearmanovo $\acs{var-rho}$ jako robustní neparametrická alternativa rezistentní vůči odlehlým hodnotám, a~$\acs{var-r}_{2024}$ = Pearsonovo $\acs{var-r}$ pro průřezová data z~roku 2024.
Síla korelací je hodnocena dle Evansovy (1996) škály: $|\acs{var-r}| < 0{,}20$ velmi slabá; $0{,}20$--$0{,}39$ slabá; $0{,}40$--$0{,}59$ střední; $0{,}60$--$0{,}79$ silná; $\geq 0{,}80$ velmi silná.

Výsledky jsou interpretovány po jednotlivých párech:

\begin{itemize}
  \item \textbf{Pokrytí \acs{KS} vs.~pracovní doba} ($\acs{var-r} = -0{,}483$; $\acs{var-rho} = -0{,}525$):
    panelová Pearsonova korelace svděčí o~\emph{střední záporné} asociaci --- vyšší pokrytí \ac{KS} je systematicky spojeno s~kratší průměrnou pracovní dobou.
    Spearmanovo $\acs{var-rho}$ je v~absolutní hodnotě vyšší než Pearsonovo $\acs{var-r}$, což indikuje robustní, přibližně monotónní vztah bez dominantních outlierů posouvájících Pearsonův koeficient.
    Průřezová korelace $\acs{var-r}_{2024} = -0{,}324$ (\emph{slabá--střední záporná}) je oproti panelové hodnotě slabematch --- v~aktuálním průřezu roku 2024 se tedy vztah projevuje méně výrazně, pravděpodobně vlivem strukturní konvergence pracovní doby v~rámci EU.

  \item \textbf{Pokrytí \acs{KS} vs.~\ac{HDP}/ob.~(\ac{PPS})} ($\acs{var-r} = 0{,}566$; $\acs{var-rho} = 0{,}472$):
    \emph{střední až silná kladná} asociace na celém panelu; Pearsonovo $\acs{var-r}$ je o~0,094 vyšší než Spearmanovo $\acs{var-rho}$, což naznačuje, že panelový výsledek je citělný na několik hodnotnějších pozorovaní (nejpravděpodobněji severskemi zeměmi s~extrémně vysokým pokrytím i~\ac{HDP}).
    Vyřazení Irska posillo panelový signál (Irsko kombinovalo nízké pokrytí s~extrémním \ac{HDP}, což dříve tlumilo kladnou korelaci); průřezová hodnota $\acs{var-r}_{2024} = 0{,}448$ (\emph{střední kladná}) v~roce 2024 zůstává statisticky středně silná.

  \item \textbf{Pokrytí \acs{KS} vs.~Gini} ($\acs{var-r} = -0{,}303$; $\acs{var-rho} = -0{,}238$):
    panelová korelace je \emph{slabá záporná}; divergence mezi Pearsonem a~Spearmanem se zúžila na 0,065 (dříve 0,159 při zahrnutí Irska), což potvrzuje, že Irsko bylo strukturálním outliernm narušujícím tuto korelaci.
    Průřezová hodnota $\acs{var-r}_{2024} = +0{,}481$ je kladná, což na první pohled odporuje teorii mzdové komprese; vysvětlení je institucionální: v~průřezu roku 2024 mají státy s~erga omnes pokrytím (Francie, Itálie, Belgie, Rakousko a~Španělsko --- všechny $\geq 90\,\%$) vyšší disponibilní Gini než řada středoevropských zemí s~nízkým pokrytím, a~to proto, že jejich Gini disponibilního příjmu je taženo nahoru vysokými kapitálovými příjmy a~odměnami ve finančním sektoru, nikoliv absencí mzdové komprese.
    Panelový koeficient ($-0{,}303$) sledující \emph{změny} v~čase v~rámci zemí je proto informačně hodnotnější: dokumentuje, že kdykoli pokrytí \ac{KS} v~dané zemi kleslo, Gini se tendovalo zvýšit.
    Příčinný mechanismus KV~$\to$~Gini tedy operuje zejména přes mzdovou distribuci (primární příjmy), nikoliv přes disponibilní příjmy ovlivněné redistribucí.

  \item \textbf{Pokrytí \acs{KS} vs.~příjmem/\ac{HDP}} ($\acs{var-r} = 0{,}210$; $\acs{var-rho} = 0{,}177$):
    \emph{slabá kladná} asociace, přibližně konzistentní v~obou metodách (rozdíl 0,033).
    Relativně nízká hodnota je z~velké části artefaktem konstrukce ukazatele: \ac{HDP}/ob.~je sám o~sobě pozitivně korelovan s~\ac{KV} pokrytím ($\acs{var-r} = 0{,}566$), takže ve jmenovateli podílu čistý~příjem/\ac{HDP} oba členové rostou společně s~pokrytím a~efekt se vzájemně tlumí.
    Tento výsledek tedy nesvědčí o~tom, že \ac{KV} instituce nezvedají mzdy, ale o~tom, že tento konkrétní relativní ukazatel není citlivým testem dané hypotézy --- absolutní příjmová míra (předchozí pár) zůstává spolehlivějším indikátorem.

  \item \textbf{Odborová organizovanost vs.~Gini} ($\acs{var-r} = -0{,}458$; $\acs{var-rho} = -0{,}462$):
    tento pár je statisticky nejrobustjším v~celé tabulce: Pearsonovo $\acs{var-r}$ a~Spearmanovo $\acs{var-rho}$ jsou prakticky totožné (rozdíl 0,004), což svědčí o~přibližně lineárním vztahu bez systematického vlivu odlehjících hodnot.
    Průřezová korelace $\acs{var-r}_{2024} = -0{,}627$ dosáží \emph{silné záporné} asociace (Evans: $0{,}60$--$0{,}79$), výrazně posilující celkový závěr: v~roce 2024 je odborová organizovanost nejsilnějším institucionálním prediktorem nízkého Gini ze všech testovaných párů.
    Shoda panelových koeficientů a~jejich posesílení v~průřezu rovněž zpětně potvrzuje, že kladná průřezová hodnota u~páru \ac{KV}~vs.~Gini je specifická pro pokrytí (kde strukturní outlieré jako Francie deformují průřezový vztah), nikoliv obecná vlastnost vztahu institucí \ac{KV} k~nerovnosti.
\end{itemize}

Žádná z~výše uvedených korelací nefalzifikuje teorii mzdové komprese prostřednictvím kolektivního vyjednávání.
Klíčovým argumentem je, že jednoduchou bivariátní korelací na heterogenním průřezu 26~zemí lze jen obtížně izolovat marginální příspěvek \ac{KV} pokrytí --- vzdor tomu jsou všechny \emph{panelové} korelace ve správném předpokládaném směru bez jediné výjimky (průřezová hodnota pro pár KV vs.~Gini je kladná, což však odraží redistribuce a~kapitálové příjmy, nikoliv mzdovou kompresi --- viz výše).
Kauzální identifikaci přináší panelová analýza s~fixními efekty a~přirozené experimenty: metaanalytická literatura (např.~Jaumotte \& Osorio Buitron 2015; Card et al.~2004) konzistentně dokládá statisticky i~věcně významný efekt \ac{KV} institucí na mzdovou kompresi.
Korelace v~tabulce tedy plní funkci deskriptivního podkladu motivujícího institucionální případ\-ovou studii ČR --- nikoliv jako izolovaný kauzální důkaz, ale jako konzistentní mezinárodní vzorec, v~němž ČR se svým pokrytím $\approx 32\,\%$ systematicky vykazuje horšíu výsledky než srovnatelné ekonomiky s~vyšším pokrytím \ac{KS}.

Tabulka \ref{tab:korelace_tabulka} shrnuje výsledky korelační analýzy provedené na párovém panelu zemí OECD/\ac{EU} za dostupná roky od roku 2004.
Zdrojem dat pro pokrytí kolektivními smlouvami je databáze OECD/AIAS ICTWSS (upravená míra \textit{AdjCov} pro většinu zemí \acs{geo-EU27}; pro Německo a~Slovensko průzkumová míra ERB, viz popis dat); pro odborovou organizovanost tatáž databáze (TUD); výstupní proměnné pocházejí z~Eurostatu (lfsa\_ewhun2, nama\_10\_pc, ilc\_di12, earn\_nt\_net).
Z~analýzy je vyřazeno Irsko: jeho \ac{HDP}/ob.~v~\acs{PPS} $\approx 237$ (\acs{geo-EU27}\,=\,100) je zkresleno přesunem zisků nadnárodních korporací se sídlem v~Irsku (transfer pricing, relokace duševního vlastnictví) --- nereprezentuje skutečnou životní úroveň ani podmínky trhu práce.
Estonsko a~Lotyšsko chybí kvůli neúplným hodnotám pokrytí \acs{KS} v~ICTWSS~\textit{AdjCov}; pro Lucembursko jsou data pokrytí dostupná pouze před rokem 2004.
Každý řádek reportuje: počet párových pozorování $\acs{var-n}$ (geo$\times$rok po párování datových sad), Pearsonovo $\acs{var-r}$ na celém panelu (míra lineární asociace), Spearmanovo $\acs{var-rho}$ jako robustní neparametrická alternativa rezistentní vůči odlehlým hodnotám, a~$\acs{var-r}_{2024}$ = Pearsonovo $\acs{var-r}$ pro průřezová data z~roku 2024.
Síla korelací je hodnocena dle Evansovy (1996) škály: $|\acs{var-r}| < 0{,}20$ velmi slabá; $0{,}20$--$0{,}39$ slabá; $0{,}40$--$0{,}59$ střední; $0{,}60$--$0{,}79$ silná; $\geq 0{,}80$ velmi silná.

Výsledky jsou interpretovány po jednotlivých párech:

\begin{itemize}
  \item \textbf{Pokrytí \acs{KS} vs.~pracovní doba} ($\acs{var-r} = -0{,}483$; $\acs{var-rho} = -0{,}525$):
    panelová Pearsonova korelace svděčí o~\emph{střední záporné} asociaci --- vyšší pokrytí \ac{KS} je systematicky spojeno s~kratší průměrnou pracovní dobou.
    Spearmanovo $\acs{var-rho}$ je v~absolutní hodnotě vyšší než Pearsonovo $\acs{var-r}$, což indikuje robustní, přibližně monotónní vztah bez dominantních outlierů posouvájících Pearsonův koeficient.
    Průřezová korelace $\acs{var-r}_{2024} = -0{,}324$ (\emph{slabá--střední záporná}) je oproti panelové hodnotě slabematch --- v~aktuálním průřezu roku 2024 se tedy vztah projevuje méně výrazně, pravděpodobně vlivem strukturní konvergence pracovní doby v~rámci EU.

  \item \textbf{Pokrytí \acs{KS} vs.~\ac{HDP}/ob.~(\ac{PPS})} ($\acs{var-r} = 0{,}566$; $\acs{var-rho} = 0{,}472$):
    \emph{střední až silná kladná} asociace na celém panelu; Pearsonovo $\acs{var-r}$ je o~0,094 vyšší než Spearmanovo $\acs{var-rho}$, což naznačuje, že panelový výsledek je citělný na několik hodnotnějších pozorovaní (nejpravděpodobněji severskemi zeměmi s~extrémně vysokým pokrytím i~\ac{HDP}).
    Vyřazení Irska posillo panelový signál (Irsko kombinovalo nízké pokrytí s~extrémním \ac{HDP}, což dříve tlumilo kladnou korelaci); průřezová hodnota $\acs{var-r}_{2024} = 0{,}448$ (\emph{střední kladná}) v~roce 2024 zůstává statisticky středně silná.

  \item \textbf{Pokrytí \acs{KS} vs.~Gini} ($\acs{var-r} = -0{,}303$; $\acs{var-rho} = -0{,}238$):
    panelová korelace je \emph{slabá záporná}; divergence mezi Pearsonem a~Spearmanem se zúžila na 0,065 (dříve 0,159 při zahrnutí Irska), což potvrzuje, že Irsko bylo strukturálním outliernm narušujícím tuto korelaci.
    Průřezová hodnota $\acs{var-r}_{2024} = +0{,}481$ je kladná, což na první pohled odporuje teorii mzdové komprese; vysvětlení je institucionální: v~průřezu roku 2024 mají státy s~erga omnes pokrytím (Francie, Itálie, Belgie, Rakousko a~Španělsko --- všechny $\geq 90\,\%$) vyšší disponibilní Gini než řada středoevropských zemí s~nízkým pokrytím, a~to proto, že jejich Gini disponibilního příjmu je taženo nahoru vysokými kapitálovými příjmy a~odměnami ve finančním sektoru, nikoliv absencí mzdové komprese.
    Panelový koeficient ($-0{,}303$) sledující \emph{změny} v~čase v~rámci zemí je proto informačně hodnotnější: dokumentuje, že kdykoli pokrytí \ac{KS} v~dané zemi kleslo, Gini se tendovalo zvýšit.
    Příčinný mechanismus KV~$\to$~Gini tedy operuje zejména přes mzdovou distribuci (primární příjmy), nikoliv přes disponibilní příjmy ovlivněné redistribucí.

  \item \textbf{Pokrytí \acs{KS} vs.~příjmem/\ac{HDP}} ($\acs{var-r} = 0{,}210$; $\acs{var-rho} = 0{,}177$):
    \emph{slabá kladná} asociace, přibližně konzistentní v~obou metodách (rozdíl 0,033).
    Relativně nízká hodnota je z~velké části artefaktem konstrukce ukazatele: \ac{HDP}/ob.~je sám o~sobě pozitivně korelovan s~\ac{KV} pokrytím ($\acs{var-r} = 0{,}566$), takže ve jmenovateli podílu čistý~příjem/\ac{HDP} oba členové rostou společně s~pokrytím a~efekt se vzájemně tlumí.
    Tento výsledek tedy nesvědčí o~tom, že \ac{KV} instituce nezvedají mzdy, ale o~tom, že tento konkrétní relativní ukazatel není citlivým testem dané hypotézy --- absolutní příjmová míra (předchozí pár) zůstává spolehlivějším indikátorem.

  \item \textbf{Odborová organizovanost vs.~Gini} ($\acs{var-r} = -0{,}458$; $\acs{var-rho} = -0{,}462$):
    tento pár je statisticky nejrobustjším v~celé tabulce: Pearsonovo $\acs{var-r}$ a~Spearmanovo $\acs{var-rho}$ jsou prakticky totožné (rozdíl 0,004), což svědčí o~přibližně lineárním vztahu bez systematického vlivu odlehjících hodnot.
    Průřezová korelace $\acs{var-r}_{2024} = -0{,}627$ dosáží \emph{silné záporné} asociace (Evans: $0{,}60$--$0{,}79$), výrazně posilující celkový závěr: v~roce 2024 je odborová organizovanost nejsilnějším institucionálním prediktorem nízkého Gini ze všech testovaných párů.
    Shoda panelových koeficientů a~jejich posesílení v~průřezu rovněž zpětně potvrzuje, že kladná průřezová hodnota u~páru \ac{KV}~vs.~Gini je specifická pro pokrytí (kde strukturní outlieré jako Francie deformují průřezový vztah), nikoliv obecná vlastnost vztahu institucí \ac{KV} k~nerovnosti.
\end{itemize}

Žádná z~výše uvedených korelací nefalzifikuje teorii mzdové komprese prostřednictvím kolektivního vyjednávání.
Klíčovým argumentem je, že jednoduchou bivariátní korelací na heterogenním průřezu 26~zemí lze jen obtížně izolovat marginální příspěvek \ac{KV} pokrytí --- vzdor tomu jsou všechny \emph{panelové} korelace ve správném předpokládaném směru bez jediné výjimky (průřezová hodnota pro pár KV vs.~Gini je kladná, což však odraží redistribuce a~kapitálové příjmy, nikoliv mzdovou kompresi --- viz výše).
Kauzální identifikaci přináší panelová analýza s~fixními efekty a~přirozené experimenty: metaanalytická literatura (např.~Jaumotte \& Osorio Buitron 2015; Card et al.~2004) konzistentně dokládá statisticky i~věcně významný efekt \ac{KV} institucí na mzdovou kompresi.
Korelace v~tabulce tedy plní funkci deskriptivního podkladu motivujícího institucionální případ\-ovou studii ČR --- nikoliv jako izolovaný kauzální důkaz, ale jako konzistentní mezinárodní vzorec, v~němž ČR se svým pokrytím $\approx 32\,\%$ systematicky vykazuje horšíu výsledky než srovnatelné ekonomiky s~vyšším pokrytím \ac{KS}.

Tabulka \ref{tab:korelace_tabulka} shrnuje výsledky korelační analýzy provedené na párovém panelu zemí OECD/\ac{EU} za dostupná roky od roku 2004.
Zdrojem dat pro pokrytí kolektivními smlouvami je databáze OECD/AIAS ICTWSS (upravená míra \textit{AdjCov} pro většinu zemí \acs{geo-EU27}; pro Německo a~Slovensko průzkumová míra ERB, viz popis dat); pro odborovou organizovanost tatáž databáze (TUD); výstupní proměnné pocházejí z~Eurostatu (lfsa\_ewhun2, nama\_10\_pc, ilc\_di12, earn\_nt\_net).
Z~analýzy je vyřazeno Irsko: jeho \ac{HDP}/ob.~v~\acs{PPS} $\approx 237$ (\acs{geo-EU27}\,=\,100) je zkresleno přesunem zisků nadnárodních korporací se sídlem v~Irsku (transfer pricing, relokace duševního vlastnictví) --- nereprezentuje skutečnou životní úroveň ani podmínky trhu práce.
Estonsko a~Lotyšsko chybí kvůli neúplným hodnotám pokrytí \acs{KS} v~ICTWSS~\textit{AdjCov}; pro Lucembursko jsou data pokrytí dostupná pouze před rokem 2004.
Každý řádek reportuje: počet párových pozorování $\acs{var-n}$ (geo$\times$rok po párování datových sad), Pearsonovo $\acs{var-r}$ na celém panelu (míra lineární asociace), Spearmanovo $\acs{var-rho}$ jako robustní neparametrická alternativa rezistentní vůči odlehlým hodnotám, a~$\acs{var-r}_{2024}$ = Pearsonovo $\acs{var-r}$ pro průřezová data z~roku 2024.
Síla korelací je hodnocena dle Evansovy (1996) škály: $|\acs{var-r}| < 0{,}20$ velmi slabá; $0{,}20$--$0{,}39$ slabá; $0{,}40$--$0{,}59$ střední; $0{,}60$--$0{,}79$ silná; $\geq 0{,}80$ velmi silná.

Výsledky jsou interpretovány po jednotlivých párech:

\begin{itemize}
  \item \textbf{Pokrytí \acs{KS} vs.~pracovní doba} ($\acs{var-r} = -0{,}483$; $\acs{var-rho} = -0{,}525$):
    panelová Pearsonova korelace svděčí o~\emph{střední záporné} asociaci --- vyšší pokrytí \ac{KS} je systematicky spojeno s~kratší průměrnou pracovní dobou.
    Spearmanovo $\acs{var-rho}$ je v~absolutní hodnotě vyšší než Pearsonovo $\acs{var-r}$, což indikuje robustní, přibližně monotónní vztah bez dominantních outlierů posouvájících Pearsonův koeficient.
    Průřezová korelace $\acs{var-r}_{2024} = -0{,}324$ (\emph{slabá--střední záporná}) je oproti panelové hodnotě slabematch --- v~aktuálním průřezu roku 2024 se tedy vztah projevuje méně výrazně, pravděpodobně vlivem strukturní konvergence pracovní doby v~rámci EU.

  \item \textbf{Pokrytí \acs{KS} vs.~\ac{HDP}/ob.~(\ac{PPS})} ($\acs{var-r} = 0{,}566$; $\acs{var-rho} = 0{,}472$):
    \emph{střední až silná kladná} asociace na celém panelu; Pearsonovo $\acs{var-r}$ je o~0,094 vyšší než Spearmanovo $\acs{var-rho}$, což naznačuje, že panelový výsledek je citělný na několik hodnotnějších pozorovaní (nejpravděpodobněji severskemi zeměmi s~extrémně vysokým pokrytím i~\ac{HDP}).
    Vyřazení Irska posillo panelový signál (Irsko kombinovalo nízké pokrytí s~extrémním \ac{HDP}, což dříve tlumilo kladnou korelaci); průřezová hodnota $\acs{var-r}_{2024} = 0{,}448$ (\emph{střední kladná}) v~roce 2024 zůstává statisticky středně silná.

  \item \textbf{Pokrytí \acs{KS} vs.~Gini} ($\acs{var-r} = -0{,}303$; $\acs{var-rho} = -0{,}238$):
    panelová korelace je \emph{slabá záporná}; divergence mezi Pearsonem a~Spearmanem se zúžila na 0,065 (dříve 0,159 při zahrnutí Irska), což potvrzuje, že Irsko bylo strukturálním outliernm narušujícím tuto korelaci.
    Průřezová hodnota $\acs{var-r}_{2024} = +0{,}481$ je kladná, což na první pohled odporuje teorii mzdové komprese; vysvětlení je institucionální: v~průřezu roku 2024 mají státy s~erga omnes pokrytím (Francie, Itálie, Belgie, Rakousko a~Španělsko --- všechny $\geq 90\,\%$) vyšší disponibilní Gini než řada středoevropských zemí s~nízkým pokrytím, a~to proto, že jejich Gini disponibilního příjmu je taženo nahoru vysokými kapitálovými příjmy a~odměnami ve finančním sektoru, nikoliv absencí mzdové komprese.
    Panelový koeficient ($-0{,}303$) sledující \emph{změny} v~čase v~rámci zemí je proto informačně hodnotnější: dokumentuje, že kdykoli pokrytí \ac{KS} v~dané zemi kleslo, Gini se tendovalo zvýšit.
    Příčinný mechanismus KV~$\to$~Gini tedy operuje zejména přes mzdovou distribuci (primární příjmy), nikoliv přes disponibilní příjmy ovlivněné redistribucí.

  \item \textbf{Pokrytí \acs{KS} vs.~příjmem/\ac{HDP}} ($\acs{var-r} = 0{,}210$; $\acs{var-rho} = 0{,}177$):
    \emph{slabá kladná} asociace, přibližně konzistentní v~obou metodách (rozdíl 0,033).
    Relativně nízká hodnota je z~velké části artefaktem konstrukce ukazatele: \ac{HDP}/ob.~je sám o~sobě pozitivně korelovan s~\ac{KV} pokrytím ($\acs{var-r} = 0{,}566$), takže ve jmenovateli podílu čistý~příjem/\ac{HDP} oba členové rostou společně s~pokrytím a~efekt se vzájemně tlumí.
    Tento výsledek tedy nesvědčí o~tom, že \ac{KV} instituce nezvedají mzdy, ale o~tom, že tento konkrétní relativní ukazatel není citlivým testem dané hypotézy --- absolutní příjmová míra (předchozí pár) zůstává spolehlivějším indikátorem.

  \item \textbf{Odborová organizovanost vs.~Gini} ($\acs{var-r} = -0{,}458$; $\acs{var-rho} = -0{,}462$):
    tento pár je statisticky nejrobustjším v~celé tabulce: Pearsonovo $\acs{var-r}$ a~Spearmanovo $\acs{var-rho}$ jsou prakticky totožné (rozdíl 0,004), což svědčí o~přibližně lineárním vztahu bez systematického vlivu odlehjících hodnot.
    Průřezová korelace $\acs{var-r}_{2024} = -0{,}627$ dosáží \emph{silné záporné} asociace (Evans: $0{,}60$--$0{,}79$), výrazně posilující celkový závěr: v~roce 2024 je odborová organizovanost nejsilnějším institucionálním prediktorem nízkého Gini ze všech testovaných párů.
    Shoda panelových koeficientů a~jejich posesílení v~průřezu rovněž zpětně potvrzuje, že kladná průřezová hodnota u~páru \ac{KV}~vs.~Gini je specifická pro pokrytí (kde strukturní outlieré jako Francie deformují průřezový vztah), nikoliv obecná vlastnost vztahu institucí \ac{KV} k~nerovnosti.
\end{itemize}

Žádná z~výše uvedených korelací nefalzifikuje teorii mzdové komprese prostřednictvím kolektivního vyjednávání.
Klíčovým argumentem je, že jednoduchou bivariátní korelací na heterogenním průřezu 26~zemí lze jen obtížně izolovat marginální příspěvek \ac{KV} pokrytí --- vzdor tomu jsou všechny \emph{panelové} korelace ve správném předpokládaném směru bez jediné výjimky (průřezová hodnota pro pár KV vs.~Gini je kladná, což však odraží redistribuce a~kapitálové příjmy, nikoliv mzdovou kompresi --- viz výše).
Kauzální identifikaci přináší panelová analýza s~fixními efekty a~přirozené experimenty: metaanalytická literatura (např.~Jaumotte \& Osorio Buitron 2015; Card et al.~2004) konzistentně dokládá statisticky i~věcně významný efekt \ac{KV} institucí na mzdovou kompresi.
Korelace v~tabulce tedy plní funkci deskriptivního podkladu motivujícího institucionální případ\-ovou studii ČR --- nikoliv jako izolovaný kauzální důkaz, ale jako konzistentní mezinárodní vzorec, v~němž ČR se svým pokrytím $\approx 32\,\%$ systematicky vykazuje horšíu výsledky než srovnatelné ekonomiky s~vyšším pokrytím \ac{KS}.

\input{texparts/python/korelace_tabulka}




Výsledky jsou interpretovány po párech:
\begin{description}
  \item[Pokrytí \ac{KV} vs.\ pracovní doba]
    ($\acs{var-r} = -0{,}554$; $\acs{var-rho} = -0{,}544$; $\acs{var-Rdet} = 0{,}31$; $\acs{var-n} = 257$):
    střední záporná asociace robustně potvrzená oběma metodami.
    Pearsonovo~$\acs{var-r}$ a~Spearmanovo~$\acs{var-rho}$ jsou prakticky shodné, což svědčí o~monotónním charakteru vztahu bez dominantních outlierů posouvajících lineární složku.
    Oproti dříve používané obvyklé pracovní době (lfsa\_ewhun2) jsou zde využita data o~skutečné týdenní pracovní době (lfsa\_ewhan2), která zachycují přesčasy i~zkrácené úvazky; korelace je v~důsledku mírně silnější.
    Průřezová korelace $\acs{var-r}_{2024} = -0{,}386$ je slabší než panelová hodnota, pravděpodobně vlivem strukturní konvergence pracovní doby v~rámci EU v~posledním desetiletí (zkracování doby v~CEE, stagnace na severu).

  \item[Pokrytí \ac{KV} vs.\ gender pay gap]
    ($\acs{var-r} = -0{,}116$; $\acs{var-rho} = -0{,}186$; $\acs{var-n}_{\text{panel}} = 262$, $\acs{var-n}_{\text{geo}} = 23$):
    velmi slabá záporná asociace dle Evansovy škály.
    Slabý koeficient je pravděpodobně důsledkem heterogenity institucionálních mechanismů ovlivňujících gender pay gap (zákonné kvóty, rodičovské dávky, sektorová struktura zaměstnanosti žen) --- samotné pokrytí \ac{KS} zdaleka nepředstavuje dominantní determinant genderové mzdové disproporce.
    Směr korelace je však konzistentní s~hypotézou: koordinované \ac{KV} omezuje manažerskou diskréci při individuálním stanovování mezd a~standardizací mzdových tarifů může snižovat prostor pro genderově podmíněnou diferenciaci.
    Průřezová korelace $\acs{var-r}_{2024} = -0{,}261$ je silnější než panelová, což naznačuje, že v~posledních letech se vztah mírně zpevňuje.

  \item[Pokrytí \ac{KV} vs.\ hodinový příjem (\si{\pps\per\hour})]
    ($\acs{var-r} = 0{,}543$; $\acs{var-rho} = 0{,}496$; $\acs{var-Rdet} = 0{,}29$; $\acs{var-n}_{\text{panel}} = 238$, $\acs{var-n}_{\text{geo}} = 22$):
    střední kladná asociace dokládá, že vyšší pokrytí \ac{KS} je spojeno s~vyšším čistým hodinovým příjmem v~PPS.
    Hodinový příjem je odvozen jako podíl ročního čistého příjmu (earn\_nt\_net, standardizovaná osoba: svobodná, bezdětná, 100\,\% průměrné mzdy, po odečtení daně a~sociálních odvodů) a~skutečného ročního pracovního fondu ($\text{skutečná\_h/týden} \times 52{,}18$~týdnů/rok).
    Oproti ročnímu příjmu v~PPS zachycuje hodinová míra efektivnější koupěschopnost za odpracovanou hodinu; země s~vysokým pokrytím \ac{KS} nejen vyplácejí vyšší příjmy, ale dosahují jich při kratší pracovní době.
    Průřezová hodnota $\acs{var-r}_{2024} = 0{,}380$ zůstává slabá až střední.

  \item[Pokrytí \ac{KV} vs.\ produktivita práce (\si{\pps\per\hour})]
    ($\acs{var-r} = 0{,}627$; $\acs{var-rho} = 0{,}542$; $\acs{var-Rdet} = 0{,}39$; $\acs{var-n}_{\text{panel}} = 313$, $\acs{var-n}_{\text{geo}} = 23$):
    \emph{silná kladná} asociace --- nejsilnější ze všech čtyř párů pokrytí \ac{KV}.
    Pearsonovo~$\acs{var-r}$ je o~0,085 vyšší než Spearmanovo~$\acs{var-rho}$, divergence je pod varovným prahem.
    Průřezová hodnota $\acs{var-r}_{2024} = 0{,}510$ zůstává střední.
    Asociace je konzistentní s~teorií, podle níž koordinované mzdové vyjednávání vytváří tlak na růst produktivity: firmy pokryté odvětvovou \ac{KS} nemohou kompenzovat vyšší mzdové náklady snížením mezd, a~proto investují do kapitálové vybavenosti a~inovací.
    ČR se v~tomto panelu nachází pod průměrem EU27 jak v~produktivitě (index~$78{,}5$), tak v~pokrytí; hypotetické posílení \ac{KV} by mohlo působit jako institucionální páka pro produktivitní konvergenci.

  \item[Hustota odborů vs.\ Gini]
    ($\acs{var-r} = -0{,}458$; $\acs{var-rho} = -0{,}462$):
    statisticky nejrobustnější pár v~celé tabulce: Pearsonovo~$\acs{var-r}$ a~Spearmanovo~$\acs{var-rho}$ jsou prakticky totožné (rozdíl 0,004), což svědčí o~přibližně lineárním vztahu bez systematického vlivu odlehlých hodnot.
    Průřezová korelace $\acs{var-r}_{2024} = -0{,}627$ dosahuje \emph{silné záporné} asociace a~je nejsilnějším institucionálním prediktorem nízkého Gini ze všech testovaných párů.
    Shoda panelových koeficientů a~jejich posílení v~průřezu zpětně potvrzuje, že kladná průřezová hodnota u~páru pokrytí~vs.~Gini je specifická pro pokrytí (kde státy jako Francie deformují průřezový vztah), nikoliv obecná vlastnost vztahu \ac{KV} institucí k~nerovnosti.
\end{description}

\paragraph{Odborová hustota a příjmová nerovnost}\label{par:union_gini}

Záporná korelace hustoty odborů a~Giniho koeficientu příjmové nerovnosti (pokles Gini s~rostoucí hustotou) je jedním z~nejcitovanějších statistických argumentů v~debatě o~institucionálním uspořádání trhu práce a~je robustní i~po kontrole o~úroveň \ac{HDP} a~sociálních transferů.
ČR se řadí do skupiny zemí s~relativně nízkým Gini, ale také nízkou a~klesající odborovou hustotou --- z~čehož vyplývá, že nynější příznivá hodnota Giniho koeficientu je křehká a~závislá na jiných faktorech (zákonná minimální mzda, transfery), nikoli na silné instituci \ac{KV}.
Historická trajektorie ostatních postkomunistických zemí (PL, HU) ukazuje, že pokles hustoty bez kompenzující reformy \ac{KS} vedl k~výraznému nárůstu Gini --- tedy že ČR se pohybuje po stejné trajektorii se zpožděním zhruba 5--10 let.
Preventivním opatřením je posilování \ac{KS} pokrytí jako náhrady klesající odborové hustoty: i~v~zemích, kde jsou odbory slabé, může vysoké pokrytí \ac{KS} (erga omnes) udržovat kompresní efekt na mzdové rozdíly a~tím stabilizovat Gini (srov.~obr.~\ref{fig:eu_pokryti_kv_mapa}).

Žádná z~panelových korelací nefalzifikuje teorii mzdové komprese prostřednictvím \ac{KV} institucí.
Jednoduchou bivariátní korelací na heterogenním průřezu 24~zemí nelze izolovat marginální příspěvek pokrytí \ac{KS}; přesto jsou všechny panelové korelace ve správném předpokládaném směru bez výjimky.
Kauzální identifikaci přináší panelová analýza s~fixními efekty: metaanalytická literatura (Jaumotte \& Osorio Buitron, 2015; Card a~kol., 2004) konzistentně dokládá věcně i~statisticky významný efekt \ac{KV} institucí na mzdovou kompresi.
Korelace v~tabulce tak plní funkci deskriptivního podkladu motivujícího institucionální případovou studii ČR --- jako konzistentní mezinárodní vzorec, v~němž ČR se svým pokrytím $\approx 31\,\%$ systematicky vykazuje horší výsledky než srovnatelné ekonomiky s~vyšším pokrytím \ac{KS}.

\paragraph{Predikce vlivu na hospodářství ČR: hypotetické pokrytí~80~\%}\label{par:predikce}

Směrnice Evropského parlamentu a~Rady (EU) 2022/2041 o~přiměřených minimálních mzdách ukládá členským státům s~pokrytím \ac{KS} pod hranicí 80~\% povinnost vypracovat akční plán pro jeho zvýšení.
ČR se svým pokrytím $\approx 31\,\%$ (2023) patří k~zemím s~největším potenciálním nárůstem v~EU --- potřebný nárůst $\Delta\acs{var-x} = 80{,}0 - 30{,}7 = +49{,}3$~pp je po Litvě a~Polsku třetím největším v~EU26.

Predikovaná změna výstupních proměnných při tomto nárůstu se dle rov.~\eqref{eq:delta_pred} vypočte jako:
\begin{equation}\label{eq:model_cz}
  \Delta\acs{var-yhat}_{\text{CZ}}
    = \hat{\beta}\,(\acs{var-x}_{1} - \acs{var-x}_{\acs{idx-0}})
    = \hat{\beta} \cdot 49{,}3\ \text{pp},
\end{equation}
kde $\hat{\beta}$ jsou panelové odhady dle rov.~\eqref{eq:ols} a~$\acs{var-x}_{\acs{idx-0}} = 30{,}7\,\%$, $\acs{var-x}_{1} = 80{,}0\,\%$.
Výsledky jsou shrnuty v~tab.~\ref{tab:model_cz}.

\input{texparts/python/korelace_model_cz}

Výsledky je třeba interpretovat jako horní odhad korelovatelné části efektu, nikoliv jako kauzální predikci: $\hat{\beta}$ je odhadnut na průřezu zemí s~různými institucemi, historií a~ekonomickými strukturami a~nezachycuje zpětné vazby (general equilibrium) ani náklady přechodu.
Přesto jsou závěry konzistentní s~tezí práce a~predikovaný profil \acgen{ČR} po reformě pokrytí \ac{KV} odpovídá kombinaci charakteristik dvou současných států \ac{EU}:
\begin{itemize}
  \item nárůst čistého hodinového příjmu o~\SI{3,90}{\pps\per\hour} ($+39\,\%$) by \acacc{ČR} posunul na aktuální úroveň Itálie (\SI{13,54}{\pps\per\hour} v~roce 2024), zatímco Švédsko (\SI{17,15}{\pps\per\hour}) zůstává citelně výše;
  \item zkrácení průměrné skutečné pracovní doby o~\SI{2,1}{\hour\per\week} by \acacc{ČR} přiblížilo Švédsku (\SI{35,1}{\hour\per\week}), které kombinuje nejkratší pracovní dobu v~\ac{EU} s~produktivitou \num{107,4} bodu (EU27\,=\,100) --- hodnotou prakticky shodnou s~modelovou predikcí \num{108,3} pro \acacc{ČR};
  \item snížení \ac{GPG} o~$1{,}1$~pp je malé a~predikovaná hodnota \SI{17,4}{\percent} zůstává nad úrovní obou referenčních zemí (Itálie~\SI{5,3}{\percent}, Švédsko~\SI{11,2}{\percent}); směr efektu je však konzistentní s~hypotézou o~standardizačním efektu \ac{KV} na mzdové tarify, samotné pokrytí však nestačí -- ukazatel je v~\ac{ČR} tažen i~strukturálními faktory (odvětvová segregace, rodičovství), které \ac{KV} neřeší přímo.
\end{itemize}
Referenční profil tedy není jediná země, ale kombinace: příjmová úroveň Itálie při délce pracovní doby a~produktivitě odpovídající Švédsku. Oba tyto státy jsou přitom charakterizovány vysokým pokrytím \ac{KV} (Švédsko~\SI{87}{\percent} v~roce 2021, Itálie~\SI{100}{\percent}; dle ICTWSS \textit{AdjCov}), což odpovídá predikovanému cílovému pokrytí v~\ac{ČR}.
Reforma pokrytí \ac{KV} (zavedení mechanismu erga omnes, revize zákona o~kolektivním vyjednávání č.\,2/1991~Sb.~\cite{zakon_kv_1991}) je z~pohledu dat jednou z~nejméně nákladných dostupných politik: nevyžaduje přímé státní výdaje a~generuje distribuční efekty srovnatelné s~rozsáhlými fiskálními programy.
